\documentclass[../../main.tex]{subfiles}
\begin{document}

\section{圆周束的基本群}\label{section:ch11-圆周束的基本群}

本节我们引入圆周束这样一个概念，并计算它的基本群。

\begin{definition}[圆周束]\label{definition:ch11-wedge-of-circles}
假设 $X$ 是一个可以表示成单位圆周 $S^1$ 的有限个同胚像 $S_1,\cdots,S_n$ 之并的 Hausdorff 空间。若存在 $X$ 的一个点 $p$ 使得只要 $i\neq j$ 便有 $S_i\cap S_j=\{p\}$，则称 $X$ 是圆周 $S_1,\cdots,S_n$ 的束（wedge）。
\end{definition}

由于每一个空间 $S_i$ 是紧致的，从而它是空间 $X$ 中的一个闭集。此外，$X$ 可以嵌入平面，如果 $C_i$ 表示 $\mathbb{R}^2$ 中以 $(i,0)$ 为圆心、以 $i$ 为半径的圆周，则 $X$ 同胚于 $C_1\cup\cdots\cup C_n$。

\begin{theorem}\label{theorem:ch11-finite-wedge-fundamental-group}
设 $X$ 是圆周 $S_1,\cdots,S_n$ 的束，$p$ 是这些圆周的公共点。则 $\pi_1(X,p)$ 是一个自由群。若 $S_i$ 中的回路 $f_i$ 为 $\pi_1(S_i,p)$ 的生成元，则回路 $f_1,\cdots,f_n$ 代表 $\pi_1(X,p)$ 的一个自由生成元组。
\end{theorem}

\par
\begin{figure}[htbp]
\centering
\includegraphics[width=0.7\textwidth]{image/fig-71-1.png}
\caption{}\label{figure:ch11-fig-71-1}
\end{figure}

\begin{proof}
当 $n=1$ 时结论显然成立。我们对 $n$ 进行归纳。证明方法类似于前一节例 1 中使用的方法。

设 $X$ 是圆周 $S_1,\cdots,S_n$ 的束、$p$ 是这些圆周的公共点。对于每一个 $i$，在 $S_i$ 上选一个异于 $p$ 的点 $q_i$。设 $W_i=S_i-q_i$，并且设
\begin{align*}
U=S_1\cup W_2\cup\cdots\cup W_n\quad\text{和}\quad V=W_1\cup S_2\cup\cdots\cup S_n,
\end{align*}
则 $U\cap V=W_1\cup\cdots\cup W_n$。参见图 71.1。空间 $U$，$V$ 以及 $U\cap V$ 都是道路连通的，因为它是具有一个公共点的道路连通空间的并。


空间 $W_i$ 同胚于一个开区间，所以 $p$ 是它的一个形变收缩核。设 $F_i:W_i\times I\to W_i$ 为这个形变收缩，由这些 $F_i$ 共同定义了一个映射 $F:(U\cap V)\times I\to U\cap V$，它是一个从 $U\cap V$ 到 $p$ 的形变收缩。（为了证明 $F$ 的连续性，只需注意 $S_i$ 为 $X$ 的闭子空间，从而 $W_i=S_i-q_i$ 为 $U\cap V$ 的闭子空间，由此 $W_i\times I$ 为 $(U\cap V)\times I$ 的闭子空间。然后应用黏结引理。）由此可见 $U\cap V$ 是单连通的，从而相对于内射诱导的单同态而言，$\pi_1(X,p)$ 是群 $\pi_1(U,p)$ 和 $\pi_1(V,p)$ 的自由积。

类似的论证说明，$S_1$ 是 $U$ 的一个形变收缩核，$S_2\cup\cdots\cup S_n$ 为 $V$ 的一个形变收缩核。所以 $\pi_1(U,p)$ 是无限循环群，回路 $f_1$ 表示一个生成元。此外，由归纳假设可见 $\pi_1(V,p)$ 是一个自由群，回路 $f_2,\cdots,f_n$ 表示一个自由生成元组。根据定理 69.2，本定理证明完成。
\end{proof}

上述定理可以推广到具有一个公共点的无穷多个圆周的并 $X$ 上。这里我们要小心地关注 $X$ 的拓扑。

\begin{definition}[相通拓扑]\label{definition:ch11-coherent-topology}
设空间 $X$ 是子空间 $X_\alpha$（$\alpha\in J$）之并。$X$ 的拓扑称为是与子空间 $X_\alpha$ 相通的（coherent），倘若子集 $C$ 满足条件：对任何一个 $\alpha$，$C\cap X_\alpha$ 是 $X_\alpha$ 中的闭集，则 $C$ 为 $X$ 中的闭集。一个等价表述是：若 $X$ 中的一个子集与每一个 $X_\alpha$ 的交为 $X_\alpha$ 中的开集，则这个子集为 $X$ 中的开集。
\end{definition}

如果 $X$ 是有限多个闭子空间 $X_1,\cdots,X_n$ 的并，则 $X$ 的拓扑必定与这些子空间相通，这是因为，如果 $C\cap X_i$ 是 $X_i$ 中的闭集，则它也是 $X$ 中的闭集，并且 $C$ 为有限多个集合 $C\cap X_i$ 之并。

\begin{definition}[圆周族的束]\label{definition:ch11-infinite-wedge}
设空间 $X$ 可以表述为某些同胚于单位圆周的子空间 $S_\alpha$（$\alpha\in J$）之并，并且存在 $X$ 中的一点 $p$ 使得只要 $\alpha\neq\beta$ 便有 $S_\alpha\cap S_\beta=\{p\}$。如果 $X$ 的拓扑与诸子空间 $S_\alpha$ 相通，则称 $X$ 为圆周族 $S_\alpha$ 的束（wedge）。
\end{definition}

对于有限的情形，定义中的相通条件可以用 Hausdorff 条件替代。此时相通条件已经被蕴涵。而对于无限的情形，这种蕴涵关系不再成立，由此我们将相通条件作为定义的一部分。我们也希望有 Hausdorff 条件，然而，它已不再必要，因为它被相通条件所蕴涵。

\begin{lemma}\label{lemma:ch11-wedge-regular}
设 $X$ 为圆周 $S_\alpha$（$\alpha\in J$）的束，则 $X$ 是正规的。此外，$X$ 的任何一个紧致子空间都包含在有限多个圆周 $S_\alpha$ 的并之中。
\end{lemma}

\begin{proof}
显然，单点集是 $X$ 中的闭集。设 $A$ 和 $B$ 为 $X$ 中无交的两个闭集，$B$ 不包含点 $p$。选取 $S_\alpha$ 中的无交子集 $U_\alpha$ 和 $V_\alpha$ 使得它们在 $S_\alpha$ 中是开的，并且分别包含 $\{p\}\cup(A\cap S_\alpha)$ 和 $B\cap S_\alpha$。设 $U=\bigcup U_\alpha$ 及 $V=\bigcup V_\alpha$，则 $U$ 与 $V$ 无交。由于所有的 $U_\alpha$ 都含有 $p$，所以 $U\cap S_\alpha=U_\alpha$。由于所有的 $V_\alpha$ 都不包含 $p$，所以 $V\cap S_\alpha=V_\alpha$。因此，$U$ 和 $V$ 为 $X$ 中的开集。因而，$X$ 是正规的。

假设 $C$ 是 $X$ 的一个紧致子空间。对于每一个 $\alpha$，倘若 $C\cap(S_\alpha-p)$ 非空，我们在其中选取一点 $x_\alpha$。这时 $D=\{x_\alpha\}$ 为 $X$ 中的闭集，这是由于它与每一个子空间 $S_\alpha$ 的交为单点集或者空集。同理，$D$ 的每一个子集都是 $X$ 中的闭集。因此，$D$ 是 $X$ 中包含于 $C$ 的闭离散子空间，由于 $C$ 是极限点紧致的，所以 $D$ 一定是有限的。
\end{proof}

\begin{theorem}\label{theorem:ch11-infinite-wedge-fundamental-group}
设 $X$ 为圆周 $S_\alpha$（$\alpha\in J$）的束，$p$ 是这些圆周的公共点，则 $\pi_1(X,p)$ 为自由群。如果 $S_\alpha$ 中的回路 $f_\alpha$ 是 $\pi_1(S_\alpha,p)$ 的一个生成元，则回路 $\{f_\alpha\}$ 代表 $\pi_1(X,p)$ 的一个自由生成元组。
\end{theorem}

\begin{proof}
设 $i_\alpha:\pi_1(S_\alpha,p)\to\pi_1(X,p)$ 是由内射诱导的同态，$G_\alpha$ 是 $i_\alpha$ 的像。

如果 $f$ 是 $X$ 中以 $p$ 为基点的一条回路，则 $f$ 的像是紧致的，从而 $f$ 在某有限多个子空间 $S_\alpha$ 的并之中。进而，如果 $f$ 和 $g$ 为 $X$ 中两条道路同伦的回路，则它们在某有限多个子空间 $S_\alpha$ 的并中也是道路同伦的。

这蕴涵着群 $\{G_\alpha\}$ 生成 $\pi_1(X,p)$。因为若 $f$ 为 $X$ 中的一条回路，则存在某有限指标集使得 $f$ 在 $S_{\alpha_1}\cup\cdots\cup S_{\alpha_n}$ 中，这时，定理 71.1 蕴涵着 $[f]$ 是群 $G_{\alpha_1},\cdots,G_{\alpha_n}$ 中某些元素的乘积。类似地，$i_\beta$ 是一个单同态。因为如果 $f$ 为 $S_\beta$ 中的一条回路并且在 $X$ 中道路同伦于一条常值道路，则 $f$ 在某有限多个子空间 $S_\alpha$ 的并中道路同伦于一条常值道路，因此定理 71.1 蕴涵着 $f$ 在 $S_\beta$ 中道路同伦于一条常值道路。

最后，假设存在一个由群 $G_\alpha$ 中元素构成的表示 $\pi_1(X,p)$ 中单位元的非空的约化字
\begin{align*}
w=(g_{\alpha_1},\cdots,g_{\alpha_n}).
\end{align*}
设 $f$ 为 $X$ 中的一条回路。$w$ 表示它的道路同伦类，则 $f$ 在 $X$ 中道路同伦于一条常值道路，从而它在某有限多个子空间 $S_\alpha$ 的并中也道路同伦于一条常值道路。这与定理 71.1 矛盾。
\end{proof}

上述定理依赖于 $X$ 的拓扑与子空间族 $S_\alpha$ 相通这一事实。考虑以下例子：

\par
\begin{figure}[htbp]
\centering
\includegraphics[width=0.5\textwidth]{image/fig-71-2.png}
\caption{}\label{figure:ch11-fig-71-2}
\end{figure}

\begin{example}[无限耳环]\label{example:ch11-infinite-earring}
设 $C_n$ 为 $\mathbb{R}^2$ 中以 $(1/n,0)$ 为圆心、以 $1/n$ 为半径的圆周。设 $X$ 为 $\mathbb{R}^2$ 的子空间并且是这些圆周的并，则 $X$ 是可数无限多个圆周之并且其中任何两个都交于原点 $p$。然而，$X$ 不是这些圆周 $C_n$ 的束。为简便起见，我们称 $X$ 为一个无限耳环（infinite earring）。

可以直接验证 $X$ 的拓扑不与子空间族 $C_n$ 相通。正半 $x$ 轴与 $X$ 的交恰好含有每一个圆周 $C_n$ 的一个点，但是这个集合不是 $X$ 中的闭集。我们换一个办法来论证这个结论。对于每一个 $n$，设 $f_n$ 为表示 $\pi_1(C_n,p)$ 的生成元的 $C_n$ 中的一条回路。我们证明 $\pi_1(X,p)$ 不是以 $\{[f_n]\}$ 为自由生成元组的自由群。事实上，我们将证明这些 $[f_i]$ 并不生成群 $\pi_1(X,p)$。

考虑依以下方式定义的 $X$ 中的回路 $g$：对于每一个 $n$，在区间 $[1/(n+1),1/n]$ 上，定义 $g$ 为这个区间到 $[0,1]$ 上的正线性映射与 $f_n$ 的复合，这给出了 $g$ 在 $(0,1]$ 上的定义，此外，定义 $g(0)=p$。因为 $X$ 具有由 $\mathbb{R}^2$ 诱导的拓扑，可见 $g$ 是连续的。见图 71.2。我们指出：对任意给定的 $n$，$[g]$ 不属于由 $[f_1],\cdots,[f_n]$ 生成的 $\pi_1(X,p)$ 的子群 $G_n$。


选取 $N>n$，并考虑映射 $h:X\to C_N$，使得当 $x\in C_N$ 时 $h(x)=x$，当 $x\notin C_N$ 时 $h(x)=p$，则 $h$ 是连续的，并且诱导了一个将 $G_n$ 的每一个元素都映为单位元的同态 $h_*:\pi_1(X,p)\to\pi_1(C_N,p)$。另一方面，$h\circ g$ 为 $C_N$ 上的一条回路，它在区间 $[1/(N+1),1/N]$ 之外取常值，而在这个区间上等于这个区间到 $[0,1]$ 上的正线性映射与 $f_N$ 的复合。因此，$h_*([g])=[f_N]$，它生成 $\pi_1(C_N,p)$！于是，$[g]\notin G_n$。
\end{example}

在前面的定理中，我们计算了无限个圆周的束的基本群。为后面使用方便，我们现在证明这种空间是存在的。（这个结果在第 14 章中要用到。）

\begin{lemma}\label{lemma:ch11-wedge-existence}
对任意的指标集 $J$，存在一个空间 $X$ 使得它是圆周 $S_\alpha$（$\alpha\in J$）的束。
\end{lemma}

\begin{proof}
赋予指标集 $J$ 以离散拓扑。设 $E$ 为积空间 $S^1\times J$。取一点 $b_0\in S^1$，设 $X$ 为将 $E$ 的闭子集 $P=b_0\times J$ 黏成一点 $p$ 所得到的商空间。设 $\pi:E\to X$ 为商映射。令 $S_\alpha=\pi(S^1\times\alpha)$。我们来证明每一个 $S_\alpha$ 同胚于 $S^1$ 并且 $X$ 是这些圆周 $S_\alpha$ 的束。

如果 $C$ 是 $S^1\times\alpha$ 中的闭子集，则 $\pi(C)$ 是 $X$ 中的闭子集。如果点 $b_0\times\alpha$ 不在 $C$ 中，则 $\pi^{-1}\pi(C)=C$。否则，$\pi^{-1}\pi(C)=C\cup P$。对每一种情形，$\pi^{-1}\pi(C)$ 都是 $S^1\times J$ 中的闭集，从而 $\pi(C)$ 是 $X$ 中的闭集。

这蕴涵着 $S_\alpha$ 本身是 $X$ 中的闭集，这是因为 $S^1\times\alpha$ 是 $S^1\times J$ 中的闭集，$\pi$ 将 $S^1\times\alpha$ 同胚地映射到 $S_\alpha$ 上。设 $\pi_\alpha$ 是这个同胚。

为了证明 $X$ 的拓扑与子空间 $S_\alpha$ 相通，设 $D\subset X$ 并且假定对于每一个 $\alpha$ 有 $D\cap S_\alpha$ 是 $S_\alpha$ 中的闭集。那么
\begin{align*}
\pi^{-1}(D)\cap(S^1\times\alpha)=\pi_\alpha^{-1}(D\cap S_\alpha),
\end{align*}
由于 $\pi_\alpha$ 连续，所以 $\pi_\alpha^{-1}(D\cap S_\alpha)$ 是 $S^1\times\alpha$ 中的闭集。于是 $\pi^{-1}(D)$ 是 $S^1\times J$ 中的闭集。根据商拓扑的定义，$D$ 是 $X$ 中的闭集。
\end{proof}

\subsection*{习题}

\begin{exercise}\label{exercise:ch11-71-01}
设空间 $X$ 可以表示成 $n$ 个各自同胚于单位圆周的子空间 $S_1,\cdots,S_n$ 之并，并且存在一点 $p$ 使得当 $i\neq j$ 时 $S_i\cap S_j=\{p\}$。
\begin{enumerate}[(a)]
\item 证明 $X$ 是 Hausdorff 空间当且仅当每一个子空间 $S_i$ 为 $X$ 中的闭集。
\item 证明 $X$ 是 Hausdorff 空间当且仅当 $X$ 的拓扑与子空间 $S_i$ 是相通的。
\item 举例说明 $X$ 可以不是 Hausdorff 空间。[提示：参见第 36 节习题 5。]
\end{enumerate}
\end{exercise}

\begin{exercise}\label{exercise:ch11-71-02}
设空间 $X$ 是闭子空间 $X_1,\cdots,X_n$ 之并，并且存在一点 $p$ 使得当 $i\neq j$ 时 $X_i\cap X_j=\{p\}$。这时我们称 $X$ 为空间 $X_1,\cdots,X_n$ 的束（wedge），并记作 $X=X_1\vee\cdots\vee X_n$。证明：如果对于每一个 $i$，$p$ 为 $X_i$ 的某开集 $W_i$ 的形变收缩核，则 $\pi_1(X,p)$ 是 $\pi_1(X_i,p)$ 相对于内射诱导的单同态的外自由积。
\end{exercise}

\begin{exercise}\label{exercise:ch11-71-03}
如果 $X$ 同胚于 $S^1$ 而 $Y$ 同胚于 $S^2$，试问关于 $X\vee Y$ 的基本群有什么结论？
\end{exercise}

\begin{exercise}\label{exercise:ch11-71-04}
证明：若 $X$ 是圆周的无限束，则 $X$ 不满足第一可数性公理。
\end{exercise}

\begin{exercise}\label{exercise:ch11-71-05}
设 $S_n$ 为 $\mathbb{R}^2$ 上以 $(n,0)$ 为圆心、以 $n$ 为半径的圆周，$\mathbb{R}^2$ 的子空间 $Y$ 是这些圆周的并，$p$ 为它们的公共点。
\begin{enumerate}[(a)]
\item 证明 $Y$ 与圆周的可数无限束 $X$ 或者例 1 中的空间不同胚。
\end{enumerate}
\end{exercise}

\end{document}